\documentclass[11pt,letterpaper]{article}
\usepackage[T1]{fontenc}
\usepackage[utf8]{inputenc}
\usepackage[spanish,es-noshorthands]{babel}
\usepackage{lmodern}
\usepackage{microtype}
\usepackage[margin=2.5cm]{geometry}
\usepackage{amsmath}
\usepackage{booktabs}
\usepackage{enumitem}
\usepackage{array}
\pagestyle{empty}
\setlength{\parindent}{0pt}
\setlength{\parskip}{6pt}

\begin{document}

{\large\textbf{Econometría II · Preparación PEP 1}}\hfill 1 de noviembre de 2026

\vspace{2pt}\hrule\vspace{10pt}

{\Large\textbf{Ejercicio 2. WOE e IV}}\hfill Tiempo: 60 min

\textit{Ejercicio de práctica con datos ficticios; no proviene de una prueba del profesor.}

\vspace{4pt}
Un banco analiza la edad como variable de un modelo de \textit{scoring} de créditos de consumo. Un cliente es «malo» si cae en mora de 90 días o más durante la ventana de observación, y «bueno» en caso contrario.

\begin{center}
\begin{tabular}{l>{\raggedleft\arraybackslash}p{2cm}>{\raggedleft\arraybackslash}p{2cm}}
\toprule
\textbf{Tramo de edad} & \textbf{Buenos} & \textbf{Malos} \\
\midrule
18-24    & 480   & 290 \\
25-34    & 1.040 & 200 \\
35-44    & 880   & 250 \\
45-54    & 960   & 180 \\
55 o más & 640   & 80  \\
\bottomrule
\end{tabular}
\end{center}

Usa
\[
\text{WOE}=\ln\!\left(\frac{\%\,\text{buenos}}{\%\,\text{malos}}\right),
\]
donde cada porcentaje es la distribución sobre el total de buenos o de malos.

\begin{enumerate}[label=\textbf{\alph*)},leftmargin=*,itemsep=8pt]
  \item Calcula el WOE y el IV de cada tramo, y el IV total.
  \item Lee el IV total con los umbrales de Siddiqi. ¿La edad predice?
  \item ¿El WOE es monótono? Si no lo es, fusiona los tramos que corresponda, justifica por qué esa fusión y no otra, y recalcula la tasa de malos, el WOE y el IV total.
  \item ¿Con qué valor entra en la regresión una persona de 38 años, antes y después de la fusión?
  \item ¿Cuántos $\beta$ usa la variable si entra con WOE? ¿Y si entra con \textit{dummies}? Responde con los tramos finales.
\end{enumerate}

\vspace{6pt}
\textit{Cada resultado se entrega con su frase de interpretación.}

\end{document}
